\section{Literature Review}

% Owner-authored section. This is the current branch, TheProblem. The search
% methodology and trimming criteria below are stated as fact because the
% owner ran this search and asked for this trimming; the specific papers
% named are still leads, not checked sources, until opened individually and
% logged in AI/SOURCES.md. The mathematical framework for this project's own
% model remains undecided; nothing below commits to one.

\subsection{Search methodology}

Claude searched the literature against three research questions, stated here in full since the
trimming below depends on them:

\begin{enumerate}
  \item[RQ1.] How do you estimate the true size or activity of a population at a location when
    only a fraction is observed?
  \item[RQ2.] How do you infer venue occupancy from event data when the observed users are a
    non-uniform sample?
  \item[RQ3.] Which point-process and arrival models describe bursty, clustered arrivals, and
    how do they treat groups rather than independent individuals?
\end{enumerate}

This search returned approximately 90 candidate references.

\subsection{Trimming the candidate set}

Reading 90 papers in full is not feasible at this stage. The set was trimmed to the 8 references
that the search itself flagged as shaping the method most directly, leaving the rest at
abstract-level triage, to revisit only if a specific gap in the model requires it. Each entry
below is still a lead, not a checked source; every one needs to be opened and logged in
\texttt{AI/SOURCES.md} before any claim in this paper cites it.

\begin{itemize}
  \item Brown et al. (2014), on defining a group check-in by a time window. No confirmed
    bibliographic detail (DOI, venue, page range) has surfaced yet; the paper still needs
    locating, not just reading.
  \item Chorley et al. (2016), \textit{ICWSM}, the main prior study of Untappd check-in data.
    Full bibliographic detail is likewise not yet confirmed.
  \item Cho, Ver Steeg, and Galstyan (2014), \textit{AAAI} 28(1):269--275, on a per-venue Hawkes
    model separating self-excitation, social, and exogenous effects in check-in data.
  \item Clement, Converse, and Royle (2017), \textit{Ecological and Evolution} 7(18):7304--7310,
    on two-stage detection: whether a group is observed at all, and how many of its members are.
  \item Zhuang and Mateu (2019), \textit{JRSS A} 182(3):919--942, on separating a periodic
    seasonal background from genuine self-excitation in clustered event data.
  \item Yauck, Rivest, and Rothman (2019), \textit{JASA} 114(525):105--114, on capture-recapture
    applied to app-activation traces to estimate unique visitors to a commercial venue.
  \item Timokhin, Sadrani, and Antoniou (2020), \textit{Smart Cities} 3(3):818--841, on
    validating a popularity signal against an independent ground truth, a low-cost recipe for
    exactly the ground-truth problem this project's synthetic simulator is meant to solve by
    construction instead.
  \item Daw and Pender (2019), \textit{Queueing Systems} 91(3--4):367--401, on how batch
    arrivals inflate occupancy variance relative to individual arrivals with the same mean.
\end{itemize}
